\documentclass[11pt]{article}

\usepackage[margin=1in]{geometry}
\usepackage{amsmath,amssymb}
\usepackage{booktabs}
\usepackage{longtable}
\usepackage{enumitem}
\usepackage{array}
\usepackage[font=small,labelfont=bf]{caption}
\usepackage{hyperref}
\hypersetup{colorlinks=true, linkcolor=blue, citecolor=blue, urlcolor=blue}

\setlist[itemize]{itemsep=2pt,topsep=4pt}
\setlist[enumerate]{itemsep=3pt,topsep=4pt}
\renewcommand{\arraystretch}{1.1}

\title{\textbf{Char Dham Ropeway Project} \\[4pt]
\large An Assessment on Lives and Livelihoods of the Kedarnath Yatra Workforce}
\author{}
\date{September 2026}

\begin{document}
\maketitle

\begin{quote}
\small
\noindent\textbf{What is and is not real in this document.} Section~\ref{sec:pilot}
reports a \emph{real} field pilot: 46 workers interviewed on the Kedarnath route during the
2024 Yatra season. Those numbers are data. Section~\ref{sec:plan} sets out the main study
that the pilot led to, and it is a \emph{plan} -- the instrument is built and the analysis
code runs, but no main-round fieldwork has happened yet. Section~\ref{sec:worked} is a
\emph{worked example}: the full analysis pipeline run end to end on a synthetic dataset that
was hand-parameterised to behave like a real Kedarnath round. Every figure in
Section~\ref{sec:worked} demonstrates that the method works. None of them says anything
about actual Kedarnath households, and they must be replaced with fielded data before anyone
quotes them.
\end{quote}

\tableofcontents
\bigskip

\section{Background}

The Char Dham circuit pulls several million pilgrims into Uttarakhand's high-Himalayan
shrines every year, and Kedarnath is the hardest leg of it: a roughly 16-kilometre climb from
Gaurikund to a temple sitting at close to 3,600 metres. A dense seasonal economy has grown
along that trail to carry pilgrims up and to feed and house them on the way. Pony and mule
workers, porters who carry loads and sometimes people in a \textit{kandi}, palki and dandi
bearers, tea-stall and dhaba operators, shopkeepers, hotel and lodge staff and owners,
drivers working the road-head stretch, a small number of guides, and casual wage labourers.
Almost all of what these households earn in a year, they earn between May and October.

The Government of Uttarakhand has proposed an aerial ropeway to Kedarnath. The detailed
project report prepared for the state was for a Gaurikund alignment (Uttarakhand Tourism
Development Board, 2022); the alignment that has since gone forward runs from Sonprayag, and
is expected to cut a journey of eight hours or so on foot down to well under an hour. For
pilgrims this is an unambiguous improvement. For the people whose entire working life is the
walk itself, it is not. A ropeway of this kind does not complement pony and porter work, it
substitutes for it directly.

This report does two things. It presents what a 2024 field pilot of 46 route workers already
established about who these workers are and how exposed they are, and it sets out the larger
study designed on the back of that pilot -- a study meant to measure poverty, multidimensional
deprivation, employment deprivation and forward-looking \emph{vulnerability} in this
workforce \emph{before} the ropeway begins operating, so that there is a baseline to argue
from later.

\section{Literature Review}

The study sits first inside the sustainable-livelihoods literature. Chambers and Conway
(1992) defined a livelihood as the capabilities, assets and activities a household uses to
make a living, and called it sustainable when it can absorb and recover from stress without
eroding those capabilities and assets. DFID (1999) turned this into the five-capitals
operational framework that most applied work now uses, and Ellis (2000) developed the
diversification side of it, which matters here because diversification is precisely what the
Kedarnath workforce lacks. Later applications, Krantz (2001) and Lax and Krug for instance,
stress the same interaction between human capability, tangible assets and actual economic
activity, and make the point that hardship often comes not from an absence of assets but from
an inability to convert the assets a household has into some new source of livelihood. In
tourism settings this is sharper still, because the household's assets sit inside a wider
structure of visitor activity that the household does not control (Laeis and Lemke, 2016).

That capacity to respond is what the risk and consumption-smoothing literature examines.
Poverty-dynamics work shows households moving in and out of poverty within a single year, so
a snapshot of current poverty understates real insecurity (Dercon and Krishnan, 2000).
Hoddinott and Quisumbing (2008) survey how households smooth welfare through consumption,
labour supply and asset holdings, and evidence from rural India indicates informal
risk-sharing inside villages does provide meaningful insurance, with consumption responding
roughly symmetrically to income gains and losses (Bold and Broer, 2021). The lesson is that
existing coping capacity, and not only current income, decides how a household absorbs a
shock.

Poverty measures the condition a household is in now. Vulnerability asks about the risk of
being deprived later. The Foster--Greer--Thorbecke class separates incidence, depth and
severity of monetary poverty (Foster, Greer and Thorbecke, 1984), and for India the
Rangarajan Committee methodology gives the line; its re-implementation directly on the
2022--23 Household Consumption Expenditure Survey by Sethu, Surya and Ruthu (2024) is what
this study uses. But monetary poverty on its own misses overlapping deprivation in other
dimensions. The counting approach of Alkire and Foster (2011) identifies multidimensional
poverty systematically, Alkire and Santos (2010) apply it across 104 developing countries,
and NITI Aayog (2021) fixes the indicator set and weights for India specifically. The
distinction is not academic: on roughly 90,000 Punjab households, Azeem, Mugera and Schilizzi
(2018) find 18 per cent of the multidimensionally poor are missed altogether by the matching
monetary measure. Apablaza, Nogales and Sehnbruch (2026) extend the same counting logic to
\emph{work itself}, building a five-indicator employment-deprivation measure out of access,
compensation, security, stability and conditions, which is directly relevant for a workforce
whose problem is the quality and durability of the job rather than the absence of one.

A separate strand moves from identifying current poverty to estimating the probability of
future poverty. Chaudhuri, Jalan and Suryahadi (2002) provide the foundational
cross-sectional approach, in which observable household characteristics are used to estimate
the probability that future welfare falls below a line. This is the Vulnerability as Expected
Poverty (VEP) framework used throughout. Pritchett, Suryahadi and Sumarto (2000) develop the
companion measure and make the targeting argument that follows from it, and Suryahadi and
Sumarto (2003) supply the chronic-versus-transient decomposition that this study adopts.
Fujii (2016) reviews the whole field and is the source for two design choices made here: that
location and education are the covariates that come out significant most consistently, and
that the measurement-error perturbation check is worth running. Günther and Harttgen (2009)
separate idiosyncratic from covariate shocks, which is why the shock module distinguishes
them. Gallardo (2018) then situates VEP inside a broader taxonomy and identifies a genuinely
different family -- Vulnerability by Mean Risk -- which asks whether a household's expected
welfare, once penalised for its own risk, clears the line, and makes no distributional
assumption at all. Chiwaula, Witt and Waibel (2011) implement one version with the plain
standard deviation; Gallardo's own refinement penalises only the \emph{downside} half of the
dispersion, on the reasoning that a household is not harmed by unusually good luck. Both are
estimated and compared against VEP in Section~\ref{sec:vmr}. Empirically, vulnerability rates
routinely exceed observed poverty rates and the mix of poverty-driven and risk-driven
vulnerability varies by setting (Azeem, Mugera and Schilizzi, 2016). The closest precedent to
what is attempted here is Lyons, Kass-Hanna and Montoya Castano (2023), who build a
21-indicator multidimensional livelihood index for Syrian refugee households in Lebanon and
push it through the same three-stage FGLS procedure, linking current deprivation to coping
responses that run from savings and borrowing up to the sale of productive assets. Feeny and
McDonald (2016) do the multidimensional version for Melanesia.

The last strand concerns what happens when a structural change removes the occupations
households currently live off. Work on technological change and task composition shows that
shifts in how work is organised change the demand for particular occupational tasks (Autor,
Levy and Murnane, 2003; Lewandowski, Park and Schotte, 2020). Within that literature,
occupational skills transferability predicts smaller earnings losses after displacement
(Nawakitphaitoon and Ormiston, 2016), and Gathmann and Schönberg (2010) give the task-based
distance measure that makes transferability computable from a task grid rather than from a
wage regression. More recent work measures proximity between a worker's own profile and the
requirements of alternative occupations directly, and generates occupation recommendations
that can differ from a worker's prior job (Bächli et al., 2024). Martins-Neto, Gukovas and
Fouarge weight that distance by the regional employment share of each destination occupation,
which is the only sensible way to handle a setting where the obvious destination occupation
does not exist yet. And most directly, Neffke, Nedelkoska and Wiederhold (2024) show
displacement raises occupational switching by a factor of eleven or twelve, with earnings
consequences that depend on the specific shape of the resulting mismatch -- they sort
transitions into stayers, lateral switchers, upskillers, downskillers and reskillers, a
typology this study's task module is built to reproduce.

Taken together the literature says livelihood security here depends on four things: the
assets and capabilities the household holds, its present exposure to poverty and other
deprivation, its capacity to absorb a shock, and its ability to move if the local economy
changes shape. The pilot in Section~\ref{sec:pilot} speaks to the first, third and fourth.
The main study is designed to add the second properly.

\section{Research Gap, Research Questions and Hypotheses}

\textbf{Research gap.} Very little empirical work brings the sustainable-livelihoods,
vulnerability and occupational-adaptation literatures together to assess a population's
livelihood security \emph{in anticipation of} a specific, named, already-approved structural
change, instead of after the fact. Displacement studies are nearly all retrospective, for the
obvious reason that displacement has to happen before you can study it. Kedarnath offers the
rarer case where the shock is announced, dated and geographically precise, and the affected
workforce is still in place and countable.

\textbf{Research questions.}
\begin{enumerate}
    \item What are the existing livelihood and occupational structures of households attached to the Yatra economy?
    \item What is the extent and distribution of current monetary poverty, multidimensional deprivation and employment deprivation, and of forward-looking vulnerability in each?
    \item Which household characteristics are associated with differences in vulnerability and adaptive capacity?
    \item How transferable are the existing skills of this workforce to other work that is actually available in the region -- including but not limited to the jobs the ropeway is claimed to create?
    \item Do workers intend to stay in the region or move, and what does that intention track?
\end{enumerate}

\textbf{Hypotheses.} The analysis tests whether differences in household characteristics,
livelihood structure, assets and adaptive capacity are associated with differences in
vulnerability; and, for the transferability arm, whether task overlap with regionally
available occupations predicts which workers would be structurally displaced rather than
merely inconvenienced.

\section{Phase I: The 2024 Field Pilot}\label{sec:pilot}

\subsection{Data and sample}

A structured questionnaire was administered to pilgrimage-route workers at Kedarnath during
the 2024 Yatra season. Respondents were purposively sampled across the occupational
categories visibly present at the site, with the aim of covering the full range of
trail-dependent work rather than of being representative. The achieved sample is 46
respondents and roughly 120 variables, covering income sources and amounts, household
composition, productive assets, savings and debt, past exposure to income shocks, awareness of
the ropeway, perceived post-ropeway opportunity, and a 1--5 Likert stance on the project.

\begin{table}[h]
\centering
\caption{Pilot sample composition by occupational group ($N=46$)}
\begin{tabular}{lrr}
\toprule
Occupational group & $N$ & Share (\%) \\
\midrule
Hospitality (hotel, lodge, dhaba, shop) & 22 & 47.8 \\
Animal handlers (pony, mule, palki)     & 12 & 26.1 \\
Service (guide, priest, other)          &  6 & 13.0 \\
Labour (porter / kandi)                 &  3 &  6.5 \\
Transport (driver)                      &  3 &  6.5 \\
\midrule
Total & 46 & 100.0 \\
\bottomrule
\end{tabular}
\end{table}

Hospitality is the largest group, which reflects its dominance in the visible Kedarnath
economy. Animal handlers are second, and as the rest of this section shows, they are the group
most exposed to the ropeway.

\subsection{The economy is mono-dependent}

Yatra-derived income accounts for over 95 per cent of mean income in \emph{every} occupational
group in the pilot. This is the single most important descriptive fact in the dataset. There
is essentially no off-season fallback that would buffer a fall in trail demand, and when the
pilot asked directly what people do outside the season, ``no other work'' was the most common
single answer, at around 39 per cent. Hospitality and transport workers hold the highest
median incomes, animal handlers sit at the bottom of the distribution.

Forty-one per cent of the sample (19 of 46) hold a BPL card. Among card holders savings rates
are lower and vulnerability index values higher than among the 27 non-holders. But the card
is a poor proxy for exposure -- several holders report incomes well above the formal
threshold, and several of the most trail-dependent workers do not hold one at all. This is the
finding that drives the targeting recommendation in Section~\ref{sec:policy}.

\subsection{Ropeway stance is occupationally structured}

Of the 42 non-neutral respondents, 23 oppose the ropeway (stance 1 or 2) and 19 support it
(stance 4 or 5). The split is not random across the trail.

\begin{table}[h]
\centering
\caption{Ropeway stance by occupational group, pilot survey counts ($N=46$)}
\begin{tabular}{lrrrrrr}
\toprule
 & \multicolumn{5}{c}{Likert stance (1 = strongly oppose, 5 = strongly support)} & \\
\cmidrule(lr){2-6}
Group & 1 & 2 & 3 & 4 & 5 & Total \\
\midrule
Animal      &  9 & 0 & 0 & 0 &  3 & 12 \\
Labour      &  1 & 1 & 0 & 0 &  1 &  3 \\
Transport   &  0 & 1 & 0 & 1 &  1 &  3 \\
Service     &  1 & 1 & 1 & 0 &  3 &  6 \\
Hospitality &  6 & 3 & 3 & 1 &  9 & 22 \\
\midrule
Total & 17 & 6 & 4 & 2 & 17 & 46 \\
\bottomrule
\end{tabular}
\end{table}

Animal handlers oppose at 75 per cent (9 of 12). Hospitality workers support at 45 per cent
(10 of 22), nine of them at the top of the scale. Stance correlates negatively with the
vulnerability index ($r = -0.39$).

\subsection{Opposers have more to lose and less to fall back on}

Comparing the standardised livelihood profiles of the 23 opposers against the 19 supporters,
opposers score systematically higher on vulnerability, Yatra income share and dependency
ratio, and lower on adaptive capacity, asset index and number of income sources. The two
things move together: the workers with the most to lose are the same workers with the least
capacity to absorb the loss.

Three models were fitted, all of them on a sample small enough that they should be read as
description and not inference.

\begin{itemize}
    \item \textbf{OLS on log income} ($N=46$, HC3 robust, $R^2 = 0.183$). Diversification is the clearest predictor: each additional income source a household holds is associated with roughly 35 per cent higher total income ($p = 0.013$). Occupation group points the expected way -- animal handlers earn less than hospitality workers -- but within-group variation swamps it at this sample size.
    \item \textbf{Logit on opposition} ($N=42$ non-neutral, HC3). The only predictor reaching conventional significance is the vulnerability index, with a marginal effect of $+0.604$ (dy/dx at means, $z=2.15$, $p=0.031$; the coefficient itself is $p=0.070$). Being an animal handler adds about 20 points of opposition probability over an otherwise similar hospitality worker but does not reach significance ($p=0.297$), because three animal handlers with lower vulnerability and higher assets strongly \emph{support} the project. Yatra income share adds nothing once vulnerability is controlled, suggesting trail dependence matters mainly \emph{through} household fragility. McFadden pseudo-$R^2$ is 0.160.
    \item \textbf{Classification, leave-one-out cross-validated} (appropriate at $N<50$). Random Forest reaches 61.9 per cent accuracy against a 54.8 per cent majority-class baseline; two logit specifications reach 57.1 per cent. The most predictive features are Yatra income share, the vulnerability index and the asset index, which agrees with the logit.
\end{itemize}

Principal components tells the same story geometrically. The first component explains 46.8 per
cent of variation on its own and it is a single resilience--dependence axis: many income
sources, high adaptive capacity and low trail dependence at one end, single-activity work with
no alternatives and almost all income from the season at the other. Animal handlers and
labourers sit at the dependence end. A multiple correspondence analysis, run twice -- once in
Stata with stance as an active variable, once in Python with stance used only for colour, so
the axes are not defined by stance -- puts strong opposition in the same neighbourhood as BPL
card, no savings, no productive assets and no work history outside the Yatra economy, and
strong support next to ownership, education and off-season diversification. That is a picture
of the raw data rather than a model result, and it agrees with every model above.

\subsection{Scenario analysis on a synthetic population}

Forty-six respondents are enough to establish who these workers are. They are not enough to
simulate an income shock -- with 12 animal handlers, group-level estimates would be driven by
two or three individuals. So a synthetic population of 1,000 households was generated from the
observed sample, half by smoothed bootstrap (resampling with a Silverman-bandwidth jitter on
continuous variables, binary characteristics copied without jitter) and half by a Gaussian
copula fitted to the joint distribution, which preserves the resilience--dependence structure
the PCA found. Validation accepts it: 42 of 46 checks pass, the Frobenius norm of the
correlation difference matrix is 0.596 against a ceiling of 2.0, and all 12 categorical and
binary checks pass. The four failures are distributional mismatches in discrete-heavy
variables where the copula smooths what is effectively a step function in a small sample.

Four displacement scenarios were applied to Yatra-derived income only.

\begin{table}[h]
\centering
\caption{Mean income by occupational group across scenarios (Rs.\ lakh per year, $N=1{,}000$ synthetic households)}
\begin{tabular}{lrrrr>{\raggedright\arraybackslash}p{2.6cm}}
\toprule
Group & Baseline & Mild & Moderate & Severe & Severe + retraining \\
\midrule
Animal      & 5.22 & 3.55 ($-32\%$) & 2.09 ($-62\%$) & 1.25 ($-78\%$) & 3.66 ($-32\%$) \\
Labour      & 6.38 & 4.87 ($-24\%$) & 3.35 ($-50\%$) & 2.34 ($-66\%$) & 5.26 ($-18\%$) \\
Transport   & 7.74 & 8.36 ($+8\%$)  & 6.47 ($-18\%$) & 5.71 ($-26\%$) & 6.47 ($-18\%$) \\
Service     & 6.62 & 7.41 ($+12\%$) & 6.38 ($-4\%$)  & 6.38 ($-4\%$)  & 6.38 ($-4\%$) \\
Hospitality & 6.45 & 8.71 ($+35\%$) & 7.74 ($+20\%$) & 8.07 ($+25\%$) & 7.74 ($+20\%$) \\
\bottomrule
\end{tabular}
\end{table}

Effective changes are smaller than the shock parameters applied, because mean Yatra share is
below 1.0 for most households. Counted against a monthly poverty line of Rs.\ 3,500 (Tendulkar
Committee threshold, inflation-adjusted; Planning Commission of India, 2009), the mild and
moderate scenarios induce almost no new poverty -- zero and three households respectively --
while the severe scenario pushes 121 households below the line, and every single one of them is
an animal handler. That is about 42 per cent of the synthetic animal-handler group. Retraining
eliminates poverty induction in the moderate case entirely.

The compensation exercise asks what lump sum would restore a three-month income buffer to each
losing household, which is enough to stop distress asset sales and forced borrowing during a
transition without being long-term income replacement:
\begin{equation}
C_i = \max\!\left(0,\; 3 \times \frac{y_i^{\text{pre}} - y_i^{\text{post}}}{12}\right)
\end{equation}
Aggregated, this comes to roughly Rs.\ 3.47 crore under moderate assumptions and Rs.\ 4.52
crore under severe -- a fraction of a per cent of the project budget. Animal handlers need
Rs.\ 78,331 per household (moderate) and Rs.\ 99,209 (severe); hospitality workers gain income
in both and need nothing. The cost is dominated by two groups, which is exactly what makes
occupational targeting efficient.

\subsection{What the pilot could not do}

The pilot is convenience-sampled, has one to thirteen observations per occupational group, and
skews toward owners, so its \emph{shares} should not be quoted as population quantities. Its
categories and wording are reusable, its shares are not. Beyond sampling, three substantive
gaps motivated the main study:

\begin{enumerate}
    \item \textbf{No consumption module.} The pilot measured income, and income in a seasonal migrant economy is the noisiest possible welfare measure. Every poverty and vulnerability result above rests on income, which is why the poverty counts sit at implausible-looking levels (zero poor at baseline).
    \item \textbf{No forward-looking vulnerability.} The pilot's ``vulnerability index'' is a composite of shock history, debt and low asset cover. It is a reasonable descriptive index but it is not vulnerability in the Chaudhuri et al. sense, which is an estimated \emph{probability} of falling below a line.
    \item \textbf{No task data.} Skill transferability was discussed but nothing was collected that could measure it. The scenario ``retraining'' parameter is an assumption, not an estimate.
\end{enumerate}

\section{Phase II: The Main Study -- Forward Plan}\label{sec:plan}

\subsection{Design}

Ex-ante, cross-sectional, one round. The study documents the workforce's livelihood, poverty
and vulnerability position before the ropeway operates, as a baseline that a later round can be
compared against. It does not attempt to estimate a causal effect of a ropeway that is not yet
running.

Sampling is quota-based against the fourteen occupational categories, not random, with quotas
recalibrated against a real enumeration of the route before fieldwork starts rather than
carried over from the pilot's shares. Field constraints are four enumerators, about ten days,
and one visit per respondent, which puts the realistic achieved sample in the range of 250 to
320. Priests appear on the route but are outside the scope and are not sampled. Guides are
included at a deliberately small share since most guides work outside the corridor itself.

\subsection{The instrument}

The fielded instrument is generated from a single machine-readable dictionary, which is what
guarantees the questionnaire, the Hindi translation, the data-entry form, the Stata variable
labels and the analysis code cannot drift apart. It currently holds 352 variables: 210 asked
of the respondent, 129 constructed afterwards in Stata, 8 items of paradata, and 5 flags that
exist only in the synthetic worked example. Thirteen modules:

\begin{longtable}{@{}cl@{\hspace{1.2em}}p{8.4cm}@{}}
\toprule
& Module & What it establishes \\
\midrule
\endhead
A & Respondent and household & Size, composition, headship, schooling in completed years, language of origin \\
B & Work and work history & Principal occupation on a 14-category list anchored to NCO-2015, plus free text; employment status on the PLFS four-way split; other activities held alongside \\
C & Monthly calendar & Activity and earnings for each of the twelve months, with an annual-total fallback for respondents who cannot recall month by month; remittances \\
D & Migration & Origin, usual place of residence and whether it is rural or urban, who brought the worker to the route \\
E & Household consumption & Item groups and recall periods matched to HCES, each asked twice -- a usual Yatra-season month and a usual off-season month \\
F & Housing, amenities, assets & Walls, roof, floor, water source and fetch time, sanitation type, cooking fuel, NITI's asset list, productive assets as separate binaries \\
G & Finance & Bank and post-office accounts, borrowing by lender type, outstanding debt, interest as rupees per hundred per month, insurance counts, schemes \\
H & Health & Recent morbidity with its cost and coping response, access barriers, child mortality and maternal-care indicators \\
I & Food security and shocks & The five reduced Coping Strategy Index items in WFP's own wording, asked for each season; shocks and coping as multi-select \\
J & Ropeway & Stance, and an open response recorded verbatim \\
K & Job quality & Apablaza et al.'s own items and their skip logic, feeding the five-indicator employment-deprivation score \\
L & Tasks and skills & A short three-answer task grid: done regularly now, done before elsewhere, never done \\
\bottomrule
\end{longtable}

Three design decisions are worth stating because they are not obvious and they were not free.

\emph{Everything seasonal is asked twice.} Fieldwork happens during the Yatra season. A single
``last 30 days'' consumption question would therefore capture on-site spending only and would
treat five months of worksite living as if it held all year. Consumption, remittances and the
coping-strategy items are each asked for a usual in-season month and a usual off-season month
and weighted into an annual figure. This is not a refinement; when it was introduced it moved
per-capita consumption by several hundred rupees and the headline poverty rate by more than
twenty points.

\emph{The household may be split.} During the season a large share of this workforce is living
at the worksite while the rest of the household is at home somewhere else. Questions that begin
``your household'' are unanswerable in that situation, so the in-season items ask about ``you
and anyone staying with you here'' and housing items are explicitly about the usual home rather
than about anything the enumerator can see.

\emph{The poverty line is per respondent.} India sets separate rural and urban lines. Because
the sample is majority-migrant with home locations scattered across several states and into
Nepal, the line is applied according to the respondent's own usual place of residence --
Rs.\ 2,515 rural and Rs.\ 3,639 urban (Sethu, Surya and Ruthu, 2024). One national line is
applied to everybody including non-Indian workers, on the same principle Lyons et al.\ (2023)
use for Syrian refugees in Lebanon: vulnerability is measured against the economy the worker is
actually embedded in, not the one they came from.

\subsection{Measures}

\paragraph{Monetary poverty.} Foster--Greer--Thorbecke, with $c_i$ monthly per-capita
consumption and $z_i$ the line that applies to household $i$:
\begin{equation}
P_\alpha = \frac{1}{N} \sum_{i=1}^{N} \left( \frac{z_i - c_i}{z_i} \right)^\alpha I(c_i < z_i)
\end{equation}
In words: for every household below its line, work out the shortfall as a share of the line,
then average across the whole sample. $\alpha = 0$ counts who is poor at all, $\alpha = 1$
averages how far below the poor fall, $\alpha = 2$ weights the very poorest more heavily.

\paragraph{Multidimensional deprivation.} The Alkire--Foster counting approach, run on NITI
Aayog's (2021) own indicator set and its own weights, so the result is comparable to a
published national figure rather than to a survey-specific index invented here:
\begin{equation}
s_i = \sum_j w_j \, d_{ij}, \qquad M_0 = H \times A
\end{equation}
A household is multidimensionally poor when $s_i \geq 1/3$. $H$ is the share crossing that
cut-off and $A$ is how deprived the poor are on average once they have crossed it.

NITI's National MPI has \textbf{three dimensions and twelve indicators}. Health and Education
and Standard of Living each carry one-third. Inside Health, Nutrition takes 1/6 and Child and
Adolescent Mortality and Maternal Health take 1/12 each. Inside Education, Years of Schooling
and School Attendance take 1/6 each. Standard of Living splits its third across seven
indicators at 1/21 apiece. Table~\ref{tab:niti} gives every one of them, the weight, and the
deprivation rule exactly as it is implemented in this study.

\begin{longtable}{@{}p{2.1cm}p{3.1cm}c p{6.6cm}@{}}
\caption{NITI Aayog (2021) National MPI: all three dimensions and all twelve indicators, with the deprivation rule as implemented here}\label{tab:niti}\\
\toprule
Dimension & Indicator & Weight & Deprived if \ldots \\
\midrule
\endfirsthead
\toprule
Dimension & Indicator & Weight & Deprived if \ldots \\
\midrule
\endhead
Health & Nutrition & 1/6 & \textbf{Not collected.} NITI defines this on measured height and weight. A read-aloud interview at a worksite cannot produce anthropometry. See the note below. \\
 & Child and adolescent mortality & 1/12 & Any death of a child or adolescent in the household in the last five years. \\
 & Maternal health & 1/12 & A birth in the household in the last five years where the mother did not have four antenatal visits \emph{or} did not have a skilled attendant at delivery. NITI's rule fails on either limb, not both. \\
\midrule
Education & Years of schooling & 1/6 & No member aged 10 or above has completed six years of schooling. This is a \emph{household} indicator, not the respondent's own schooling. \\
 & School attendance & 1/6 & Any child of school-going age (class 1 to 8) in the household is not attending. Households with no child in that age band are not deprived. \\
\midrule
Standard of living & Cooking fuel & 1/21 & The household cooks with dung, firewood, crop residue or shrubs, charcoal or coal. Kerosene is deliberately \emph{not} counted, because it is not on NITI's list, although the global MPI does count it. \\
 & Sanitation & 1/21 & The toilet facility is unimproved, or it is improved but shared with other households. \\
 & Drinking water & 1/21 & The source is unimproved, or it is improved but more than a 30-minute round trip away. Both limbs, not just the source. \\
 & Electricity & 1/21 & The household has no electricity. \\
 & Housing & 1/21 & The floor \emph{or} the roof \emph{or} the wall is made of inadequate material. All three limbs are asked. \\
 & Assets & 1/21 & The household owns no more than one of radio, television, telephone, computer, animal cart, bicycle, motorcycle and refrigerator, \emph{and} does not own a car, jeep or truck. \\
 & Bank account & 1/21 & No household member has a bank \emph{or post office} account. \\
\bottomrule
\end{longtable}

\noindent Two things follow from that table and both have to be carried through the whole
analysis. First, nutrition cannot be collected, so its 1/6 is redistributed proportionally over
the remaining ten indicators and the resulting score must always be reported as a
\textbf{ten-of-twelve adaptation of the National MPI}, never as the National MPI itself.
Second, several of these indicators are about the household's \emph{usual home}, which during
the Yatra season is somewhere the enumerator is not standing. They are therefore self-reported
about that home, not observed, and the questionnaire says so explicitly in the wording.

\paragraph{Employment deprivation.} Apablaza, Nogales and Sehnbruch (2026) apply the same
counting logic to work itself. Their measure has five equally weighted indicators and the
thresholds in their Table 3 are applied as written, with the departures noted below.

\begin{longtable}{@{}p{2.3cm}c p{10.3cm}@{}}
\caption{Apablaza, Nogales and Sehnbruch (2026): the five employment dimensions as implemented}\\
\toprule
Dimension & Weight & Deprived if \ldots \\
\midrule
\endfirsthead
\toprule
Dimension & Weight & Deprived if \ldots \\
\midrule
\endhead
Access & 1/5 & More than six months without work in the year \emph{and} time spent searching; or working under 20 hours a week and wanting more work. Hours are checked in both seasons, since in-season hours run 60 to 90 a week and the limb could otherwise never fire. \\
Compensation & 1/5 & Monthly work income below two-thirds of the sample median. \\
Security & 1/5 & Wage workers: no signed contract. Self-employed: the business is not registered. \\
Stability & 1/5 & Wage workers under one season in the work, self-employed under two, \emph{or} an occasional/casual job. Tenure here is counted in Yatra seasons of roughly six months, which is how Apablaza's under-six-months and under-twelve-months thresholds map onto this setting. The casual-job limb is this study's addition and not theirs. \\
Conditions & 1/5 & A workplace injury or death in the last twelve months, or neither work-related health cover nor any pension contribution. \\
\bottomrule
\end{longtable}

\noindent A worker is employment-poor at $k = 2$ of 5, which is the cut-off Apablaza et al.\ use.

\subsection{Vulnerability estimation}

Chaudhuri, Jalan and Suryahadi (2002) estimate the probability that a household's welfare
falls below the line in the near future, from a single cross-section, on two assumptions:
households with similar observables face similar average welfare and similar welfare
\emph{risk}, and both can be recovered from how actual outcomes scatter around what the
observables predict. Stage one regresses log consumption on characteristics,
\begin{equation}
\ln(c_i) = X_i'\beta + \varepsilon_i
\end{equation}
stage two regresses the squared residual on the same characteristics to get each household's
own predicted variance, and stage three re-weights the original regression by that variance
(feasible generalised least squares). Each household ends up with its own predicted mean and
its own predicted risk, and the probability follows off the normal curve:
\begin{equation}
V_i = \Phi\!\left( \frac{\ln z_i - \widehat{\mu}_i}{\widehat{\sigma}_i} \right)
\end{equation}
A household predicted to average comfortably above the line, but with a lot of uncertainty
around that average, can still carry a high probability of falling below it. Following
convention a household is called vulnerable at $V_i \geq 0.50$, and Section~\ref{sec:robust}
sweeps that threshold rather than defending it.

The same machinery is run on all three welfare measures: log consumption, the NITI-structured
deprivation score, and the Apablaza employment-deprivation score. For the two deprivation
scores the inequality flips, since a \emph{higher} score is the bad outcome. Applying
Chaudhuri's FGLS to an employment-deprivation score is, as far as this project's literature
search goes, not something any published paper has done, and it is the methodological
contribution the study is aiming at.

Covariates are organised following Azeem, Mugera and Schilizzi (2016) into adaptive capacity
(education, bank account, institutional credit, training, smartphone), sensitivity (age,
household size, employment status, years in Yatra work) and exposure (migrant status, closure
regime, shocks, Yatra income share, income seasonality, health-access tier, urban, trail
dependence). Location is in the set deliberately, on Fujii's (2016) finding that education and
location are the two covariates that come out significant most consistently across this
literature.

\paragraph{Chronic versus transient.} Suryahadi and Sumarto (2003) split on \emph{expected}
consumption against the line, not on vulnerability status: chronic poor are poor now and have
$E[c] < z$; transient poor are poor now but have $E[c] \geq z$. Vulnerability against a
threshold $\tau$ is a separate third axis, which gives a six-cell table rather than the familiar
four-cell one. An earlier version of this project's code split on vulnerability and carried
their names for a different quantity; that has been corrected.

\subsection{Occupational transferability}

This is a separate second paper and its questions in the instrument are deliberately short, but
the design is fixed. The task grid records what the worker does, not how well they rate
themselves, because a rating scale is not answerable by this population and a factual grid is.
From it, $q_{oj}$ -- the fraction of workers in occupation $o$ doing task $j$ -- gives an angular
separation between occupations exactly as in Gathmann and Schönberg (2010), and distance is one
minus that similarity. The measure is asymmetric by construction, which is the point:
Nawakitphaitoon and Ormiston's (2016) observation that $t_{ij} \neq t_{ji}$ is what lets a move
be classified as lateral, upskilling, downskilling or reskilling in Neffke, Nedelkoska and
Wiederhold's (2024) sense.

The single-destination problem -- ``can a pony handler become a ropeway operator'' -- is
dissolved by weighting distance to every other occupation by that occupation's regional
employment share, following Martins-Neto, Gukovas and Fouarge. Ropeway jobs then become one
entry in that sum, carrying the government's own claimed employment weight, so the policy claim
gets tested on its own terms instead of being assumed. The validation referees will want is the
observed-versus-random mobility test, which needs previous occupation on a broad classification
and not on the Yatra list, and that is why previous and target occupation are both free text
coded to NCO-2015 afterwards.

\subsection{Fieldwork and outputs}

Deliverables already built: a bilingual paper questionnaire, an offline data-entry form that
needs no network, enumerator scripts in English and Hindi, a question register that gives the
source and reasoning for every single item, and the full analysis chain from raw export to
final tables. Everything regenerates from the dictionary, in order, and a set of automated
checks re-verifies skip logic and variable ranges on every rebuild.

\section{Worked Example on Synthetic Data}\label{sec:worked}

Everything in this section is computed on a hand-parameterised synthetic dataset of 200
generated respondents, of which 104 survive an explicit and deliberately non-random attrition
process (weighted toward the busiest occupations, lower education, older age and migrant
status), leaving 98 with complete covariates for the regressions. The point is to show the
pipeline runs end to end and produces sensible, internally consistent output. The point is not
the numbers.

\subsection{Monetary poverty}

\begin{table}[h]
\centering
\caption{Monetary poverty, synthetic worked example ($N = 104$)}
\begin{tabular}{lr}
\toprule
Measure & Value \\
\midrule
Mean monthly per-capita consumption & Rs.\ 2,831 \\
Median & Rs.\ 2,731 \\
Applicable line: rural / urban & 86 / 18 households \\
Poverty headcount, $P_0$ & 0.404 \\
Poverty gap, $P_1$ & 0.084 \\
Squared poverty gap, $P_2$ & 0.023 \\
\bottomrule
\end{tabular}
\end{table}

Mean consumption sits only about 12 per cent above the rural line, so two in five households
fall below whichever line applies to them. Note that 26.9 per cent of respondents could not
recall twelve separate monthly earnings figures and took the annual-total fallback instead;
their income seasonality is a lower bound by construction, and the measurement-mode flag has to
be carried into any specification that uses seasonality.

\subsection{Multidimensional and employment deprivation}

\begin{table}[h]
\centering
\caption{Deprivation measures, synthetic worked example}
\begin{tabular}{lr}
\toprule
Measure & Value \\
\midrule
Multidimensional headcount, $H$ (ten of NITI's twelve indicators) & 0.077 \\
Average intensity among the poor, $A$ & 0.389 \\
Adjusted headcount ratio, $M_0$ & 0.030 \\
\midrule
Mean employment-deprivation score (Apablaza et al.) & 0.385 \\
Employment-poor headcount, $k = 2$ of 5 & 0.644 \\
\bottomrule
\end{tabular}
\end{table}

The three measures are emphatically not picking out the same households, which is the whole
reason for running all three. Multidimensional poverty is \emph{low} here -- these workers
mostly have electricity, a phone, a bank account and an improved water source at home -- while
employment deprivation is very high, because what is wrong with these jobs is their quality,
their security and their stability, and not the physical standard of living they finance.

\subsection{Vulnerability as expected poverty}

\begin{table}[h]
\centering
\caption{Suryahadi--Sumarto cells, monetary vulnerability ($N = 98$)}
\begin{tabular}{llr}
\toprule
Cell & Definition & Share \\
\midrule
A & Poor, low $E[c]$, vulnerable (chronic) & 23.5\% \\
B & Poor, high $E[c]$, vulnerable (transient) & 1.0\% \\
C & Poor, high $E[c]$, not vulnerable (transient) & 15.3\% \\
D & Non-poor, low $E[c]$, vulnerable & 19.4\% \\
E & Non-poor, high $E[c]$, vulnerable & 1.0\% \\
F & Non-poor, high $E[c]$, not vulnerable & 39.8\% \\
\bottomrule
\end{tabular}
\end{table}

Collapsing: 44.9 per cent are vulnerable at $\tau = 0.5$, 22.1 per cent are chronically poor and
15.4 per cent transiently poor. Cell D is the one policy usually misses -- one household in five
is not poor today yet has an expected consumption below its own line, which is the definition of
a household about to become poor. Multidimensional vulnerability is 3.1 per cent, and the
overlap with monetary vulnerability is almost nil (2 of the 44 monetarily vulnerable). On
synthetic data that is a mechanical consequence of a low MPI headcount and should not be read as
a finding, but it does confirm the two arms are not measuring the same thing.

The employment arm gives a vulnerability rate of 38.8 per cent, decomposing into 34.7 per cent
chronic, 4.1 per cent transient, 27.6 per cent escaped and 33.7 per cent never
employment-poor. Cross-tabulated against current monetary poverty, half of the
employment-vulnerable are currently monetary-poor against a third of the not-vulnerable --
partial overlap, not identity. A targeting rule built on either one alone would miss a large
part of the other.

\subsection{Vulnerability by Mean Risk}\label{sec:vmr}

VEP assumes log-normality and reads a probability off that curve. The mean-risk family compares
expected consumption, penalised for risk, against the line directly, in rupee levels and with no
distributional assumption. Both variants are estimated with their own two-step procedure on
levels rather than logs, so as not to smuggle the log-normality back in.

\begin{align}
\text{Chiwaula, Witt and Waibel (2011):} \quad & \text{vulnerable} \iff \widehat{\mu}_i - \widehat{\sigma}_i < z_i \\
\text{Gallardo's downside refinement:} \quad & \text{vulnerable} \iff \widehat{\mu}_i - \gamma \widehat{\sigma}_i^{-} \leq z_i, \quad \gamma \in (0,1]
\end{align}

The second counts only how far consumption tends to fall \emph{short} of its own predicted
average and ignores how far it sometimes runs above it, which separates two households with
identical means and identical overall variance if one of them is more exposed to bad outcomes
specifically.

\begin{table}[h]
\centering
\caption{Three criteria, three different pictures ($N = 98$)}
\begin{tabular}{llr}
\toprule
Criterion & Paradigm & Rate \\
\midrule
VEP (Chaudhuri, Jalan and Suryahadi, 2002) & expected poverty, log-normal & 44.9\% \\
Mean deviation (Chiwaula, Witt and Waibel, 2011) & mean risk, symmetric, levels & 84.7\% \\
Downside semi-deviation (Gallardo, 2018), $\gamma = 0.5$ & mean risk, downside only, levels & 56.1\% \\
\bottomrule
\end{tabular}
\end{table}

The choice of criterion is not a technicality. On the same 98 households three defensible
definitions of ``vulnerable'' return 45, 56 and 85 per cent. The symmetric rule is by far the
most permissive because it penalises households for upside variability they would never
experience as harm. The two criteria that are actually sensitive to which direction risk runs
agree far more closely with each other.

\subsection{Robustness}\label{sec:robust}

\begin{table}[h]
\centering
\caption{Sensitivity of the monetary vulnerability rate}
\begin{tabular}{lrrrrr}
\toprule
Vulnerability threshold $\tau$ & 0.3 & 0.4 & 0.5 & 0.6 & 0.7 \\
Rate & 68.4\% & 57.1\% & 44.9\% & 35.7\% & 24.5\% \\
\midrule
Downside criterion, $\gamma$ & 0.25 & 0.50 & 0.75 & 1.00 & \\
Rate & 44.9\% & 56.1\% & 62.2\% & 69.4\% & \\
\midrule
Share of variance treated as measurement error & 0.0 & 0.2 & 0.3 & 0.4 & 0.5 \\
Rate & 44.9\% & 44.9\% & 44.9\% & 44.9\% & 44.9\% \\
\bottomrule
\end{tabular}
\end{table}

Two of these three move a lot and one does not move at all. The threshold sweep nearly triples
the rate across its range and the $\gamma$ sweep adds twenty-five points, so any application of
either criterion should report the parameter and preferably the range, not a single point
estimate. The measurement-error perturbation, following Pritchett, Suryahadi and Sumarto (2000)
and Azeem, Mugera and Schilizzi (2016), does not move the rate at all even when half the
estimated consumption variance is treated as pure noise, which indicates classified households
sit clear of the 0.50 threshold rather than piled on top of it. That is reassuring rather than
suspicious.

\section{Discussion and Policy Implications}\label{sec:policy}

Four things come out of the pilot and the design together.

First, opposition to the ropeway is structural, not informational. The pilot's single
significant predictor of opposition is the vulnerability index, and the same cluster --
trail dependence, no off-season work, no assets, no savings, a BPL card -- shows up in the
logit, in the classifier's feature importances, in the first principal component and in the
correspondence analysis. Workers oppose the project because they have the most to lose from it
and the least capacity to absorb that loss, not because they have misunderstood it. Consultation
exercises designed around the assumption of misunderstanding will fail.

Second, prevention is cheap and sequencing is everything. A targeted scheme of Rs.\ 3.5 to 4.5
crore, on the pilot's own scenario arithmetic, restores a three-month income buffer to every
losing household and removes poverty induction entirely. That is a fraction of a per cent of the
project cost. But an income shock which arrives before the buffer does will trigger distress
asset sales, borrowing at informal rates, and out-migration -- and those are hard to reverse
with a later payment (Bryceson, 2000). Compensation and retraining have to be completed before
commercial operations begin, not after the first season.

Third, occupational registration beats BPL targeting. BPL card status overlaps with exposure in
the pilot but the overlap is partial: several card holders earn well above the threshold and
several of the most exposed workers hold no card. Eligibility keyed to pony-handler licence,
porter registration or equivalent would reach the structurally exposed regardless of their
formal poverty classification.

Fourth, retraining needs visible placements first. The retraining scenario is the best outcome
in the whole simulation, but about 30 per cent of pilot respondents already see no post-ropeway
opportunity for themselves and animal handlers are over-represented among them. A programme
announced without committed, visible placements in ropeway operations, station services or
guiding will not get uptake, and then its failure will be read as evidence that the workers did
not want to retrain.

Two further implications follow from the main study's design rather than from the pilot. The
transient group -- households that are not poor now but whose expected consumption sits below
their own line -- is invisible to any targeting rule built on current poverty status, and the
worked example puts it at around one household in five. And employment deprivation and
consumption poverty identify overlapping but distinctly different populations, so a scheme
aimed at income alone will leave the job-quality problem exactly where it found it. That
matters here because the ropeway does not reduce anyone's income to zero. It removes a
particular kind of work.

The broader point is the one the sustainable-livelihoods literature has been making since
Chambers and Conway (1992): what protects a household through a structural change is its
capacity to redeploy what it already has, not its income level on the day the change arrives.
Evidence from comparable Indian transitions, most closely the Jharkhand coal economy (Pelz et
al., 2024), suggests households do shift toward alternative skill-based livelihoods when the
opportunity is real. The task-based literature is equally clear that whether they succeed
depends on how transferable their existing skills actually are, and not on generic retraining
(Autor, Levy and Murnane, 2003; Nawakitphaitoon and Ormiston, 2016; Bächli et al., 2024).

\section{Limitations}

\begin{itemize}
    \item \textbf{The pilot is not representative.} Forty-six respondents, purposively sampled, one to thirteen per occupational group, and skewed toward owners. Its categories and question wording carry forward; its shares should not be quoted as population quantities.
    \item \textbf{The pilot measured income, not consumption.} Every poverty count in Section~\ref{sec:pilot}, including the scenario headcounts, rests on annual income in an economy where income is violently seasonal. This is the main reason the main study rebuilt the welfare measure from scratch.
    \item \textbf{Section~\ref{sec:worked} is synthetic.} It demonstrates that the pipeline works. It is not evidence about Kedarnath and must be replaced before any figure in it is quoted.
    \item \textbf{The MPI here is a ten-of-twelve adaptation.} NITI's nutrition indicator needs anthropometry, which a questionnaire cannot collect. The weight is redistributed proportionally and the result is not the National MPI, and should never be presented as comparable to it.
    \item \textbf{Consumption is not price-deflated between contexts.} In-season consumption happens at a high-altitude worksite and off-season consumption mostly happens somewhere else entirely, often in another state. The two are added together without a spatial price adjustment, because no deflator exists at the resolution this needs.
    \item \textbf{The poverty line is contested.} India has notified no official line since 2011--12. The Rangarajan re-computation on 2022--23 HCES data (Sethu, Surya and Ruthu, 2024) is used because it is the only estimate that re-runs the actual methodology on current data rather than inflating the 2011--12 figure forward, but it is a live methodological choice and should always be reported with the line stated.
    \item \textbf{The income and occupation calendar is cited to a source not yet read line by line.} Module C's monthly design follows the FAO/ILO Livelihood Assessment Tool-kit. Unlike every other source in this study, that document has not yet been checked directly against the wording actually used, and it should be before fielding.
    \item \textbf{Cross-sectional vulnerability is not prediction.} Like all VEP-style estimates, everything here assumes the relationship between characteristics and welfare risk is stable going forward. These are indicators of prospective risk, not forecasts.
\end{itemize}

\begin{thebibliography}{99}

\bibitem{alkirefoster2011} Alkire, S. and Foster, J. (2011). Counting and multidimensional poverty measurement. \textit{Journal of Public Economics}, 95(7--8):476--487.

\bibitem{alkiresantos2010} Alkire, S. and Santos, M. E. (2010). Acute Multidimensional Poverty: A New Index for Developing Countries. \textit{SSRN Electronic Journal}.

\bibitem{apablaza2026} Apablaza, M., Nogales, R., and Sehnbruch, K. (2026). The Work Dimension in Multidimensional Poverty Measurement. \textit{Social Indicators Research}, 183(2):25.

\bibitem{autor2003} Autor, D. H., Levy, F., and Murnane, R. J. (2003). The Skill Content of Recent Technological Change: An Empirical Exploration. \textit{Quarterly Journal of Economics}, 118(4):1279--1333.

\bibitem{azeem2016} Azeem, M. M., Mugera, A. W., and Schilizzi, S. (2016). Poverty and vulnerability in the Punjab, Pakistan: A multilevel analysis. \textit{Journal of Asian Economics}, 44:57--72.

\bibitem{azeem2018} Azeem, M. M., Mugera, A. W., and Schilizzi, S. (2018). Vulnerability to Multi-Dimensional Poverty: An Empirical Comparison of Alternative Measurement Approaches. \textit{The Journal of Development Studies}, 54(9):1612--1636.

\bibitem{bachli2024} B\"achli, M., Benghalem, H., Tinello, D., Aschwanden, D., Zuber, S., Kliegel, M., Pellizzari, M., and Lalive, R. (2024). Ranking occupations by their proximity to workers' profiles. \textit{Swiss Journal of Economics and Statistics}, 160(1):8.

\bibitem{bold2021} Bold, T. and Broer, T. (2021). Risk Sharing in Village Economies Revisited. \textit{Journal of the European Economic Association}, 19(6):3207--3248.

\bibitem{bryceson2000} Bryceson, D. F. (2000). Rural Africa at the Crossroads: Livelihood Practices and Policies. \textit{Natural Resource Perspectives}, 52. Overseas Development Institute.

\bibitem{chambers1992} Chambers, R. and Conway, G. (1992). Sustainable Rural Livelihoods: Practical Concepts for the 21st Century. IDS Discussion Paper 296, Institute of Development Studies, Brighton.

\bibitem{chaudhuri2002} Chaudhuri, S., Jalan, J., and Suryahadi, A. (2002). Assessing Household Vulnerability to Poverty from Cross-Sectional Data: A Methodology and Estimates from Indonesia. Discussion Paper 0102-52, Columbia University Department of Economics.

\bibitem{chiwaula2011} Chiwaula, L. S., Witt, R., and Waibel, H. (2011). An Asset-Based Approach to Vulnerability: The Case of Small-Scale Fishing Areas in Cameroon and Nigeria. \textit{The Journal of Development Studies}, 47(2):338--353.

\bibitem{dercon2000} Dercon, S. and Krishnan, P. (2000). Vulnerability, seasonality and poverty in Ethiopia. \textit{The Journal of Development Studies}, 36(6):25--53.

\bibitem{dfid1999} DFID (1999). \textit{Sustainable Livelihoods Guidance Sheets}. Department for International Development, London.

\bibitem{ellis2000} Ellis, F. (2000). \textit{Rural Livelihoods and Diversity in Developing Countries}. Oxford University Press, Oxford.

\bibitem{faoilo2009} FAO and ILO (2009). \textit{The Livelihood Assessment Tool-kit: Analysing and Responding to the Impact of Disasters on the Livelihoods of People}. Food and Agriculture Organization / International Labour Organization, Rome and Geneva.

\bibitem{feeny2016} Feeny, S. and McDonald, L. (2016). Vulnerability to Multidimensional Poverty: Findings from Households in Melanesia. \textit{The Journal of Development Studies}, 52(3):447--464.

\bibitem{foster1984} Foster, J., Greer, J., and Thorbecke, E. (1984). A Class of Decomposable Poverty Measures. \textit{Econometrica}, 52(3):761--766.

\bibitem{fujii2016} Fujii, T. (2016). Concepts and Measurement of Vulnerability to Poverty and Other Issues: A Review of Literature. Edward Elgar Publishing, Cheltenham.

\bibitem{gallardo2018} Gallardo, M. (2018). Identifying Vulnerability to Poverty: A Critical Survey. \textit{Journal of Economic Surveys}, 32(4):1074--1105.

\bibitem{gathmann2010} Gathmann, C. and Sch\"onberg, U. (2010). How General Is Human Capital? A Task-Based Approach. \textit{Journal of Labor Economics}, 28(1):1--49.

\bibitem{gunther2009} G\"unther, I. and Harttgen, K. (2009). Estimating Households' Vulnerability to Idiosyncratic and Covariate Shocks: A Novel Method Applied in Madagascar. \textit{World Development}, 37(7):1222--1234.

\bibitem{hoddinott2008} Hoddinott, J. and Quisumbing, A. R. (2008). Methods for Microeconometric Risk and Vulnerability Assessments. SSRN Working Paper 1281055.

\bibitem{krantz2001} Krantz, L. (2001). \textit{The Sustainable Livelihood Approach to Poverty Reduction: An Introduction}. Swedish International Development Cooperation Agency, Stockholm.

\bibitem{laeis2016} Laeis, G. C. M. and Lemke, S. (2016). Social entrepreneurship in tourism: applying sustainable livelihoods approaches. \textit{International Journal of Contemporary Hospitality Management}, 28(6):1076--1093.

\bibitem{laxkrug} Lax, J. and Krug, J. Livelihood Assessment: A Participatory Tool for Natural Resource Dependent Communities.

\bibitem{lewandowski2020} Lewandowski, P., Park, A., and Schotte, S. (2020). \textit{The Global Distribution of Routine and Non-Routine Work}. WIDER Working Paper 2020/75, UNU-WIDER, Helsinki.

\bibitem{lyons2023} Lyons, A. C., Kass-Hanna, J., and Montoya Castano, A. (2023). A multidimensional approach to measuring vulnerability to poverty among refugee populations. \textit{Journal of International Development}, 35(7):2014--2045.

\bibitem{martinsneto} Martins-Neto, A., Gukovas, R. M., and Fouarge, D. Forced Displacement and Occupational Mobility: A Skills-Based Approach.

\bibitem{maxwell2008} Maxwell, D. and Caldwell, R. (2008). \textit{The Coping Strategies Index: Field Methods Manual}. 2nd edition. CARE and World Food Programme.

\bibitem{mospi2024} Ministry of Statistics and Programme Implementation (2024). \textit{Household Consumption Expenditure Survey 2022--23} and \textit{2023--24}. Government of India, New Delhi.

\bibitem{nawakitphaitoon2016} Nawakitphaitoon, K. and Ormiston, R. (2016). The estimation methods of occupational skills transferability. \textit{Journal for Labour Market Research}, 49(4):317--327.

\bibitem{nco2015} National Classification of Occupations 2015, Volumes I, II-A and II-B. Directorate General of Employment and Training, Ministry of Labour and Employment, Government of India.

\bibitem{neffke2024} Neffke, F., Nedelkoska, L., and Wiederhold, S. (2024). Skill mismatch and the costs of job displacement. \textit{Research Policy}, 53(2):104933.

\bibitem{niti2021} NITI Aayog (2021). \textit{National Multidimensional Poverty Index: Baseline Report}. Government of India, New Delhi.

\bibitem{pelz2024} Pelz, S., Krumm, A., Aklin, M., Nandan, V., and Urpelainen, J. (2024). The spatial and economic footprint of the coal industry on rural livelihoods in Jharkhand, India. \textit{Energy Policy}, 186:113973.

\bibitem{planningcommission2009} Planning Commission of India (2009). \textit{Report of the Expert Group to Review the Methodology for Estimation of Poverty}. Government of India, New Delhi.

\bibitem{pritchett2000} Pritchett, L., Suryahadi, A., and Sumarto, S. (2000). Quantifying Vulnerability to Poverty: A Proposed Measure, Applied to Indonesia. World Bank Policy Research Working Paper 2437.

\bibitem{sethu2024} Sethu, C. A., Abhinav Surya, L. T., and Ruthu, C. A. (2024). Poverty in India: The Rangarajan Method and the 2022--23 Household Consumption Expenditure Survey. \textit{Review of Agrarian Studies}, 14(2):3--22.

\bibitem{suryahadi2003} Suryahadi, A. and Sumarto, S. (2003). Poverty and Vulnerability in Indonesia Before and After the Economic Crisis. \textit{Asian Economic Journal}, 17(1):45--64.

\bibitem{unhcr2025} UNHCR, UNICEF and WFP (2025). \textit{Vulnerability Assessment of Syrian Refugees in Lebanon (VASyR) 2025}. Beirut.

\bibitem{utdb2022} Uttarakhand Tourism Development Board (2022). \textit{Kedarnath Ropeway Project: Detailed Project Report}. Government of Uttarakhand, Dehradun.

\end{thebibliography}

\end{document}
